% capitulos/2metodologia.tex
\section{Metodolog\'ia}

\subsection{Dise\~{n}o Metodol\'ogico}

La presente investigaci\'on adopta un enfoque
\textbf{cuantitativo experimental}, orientado al dise\~{n}o,
implementaci\'on y validaci\'on del modelo ARGOS-11 para la
detecci\'on de personas bajo condiciones simult\'aneas de
oclusi\'on severa y baja visibilidad en im\'agenes RGB
convencionales. Seg\'un Hern\'andez Sampieri \&
Mendoza~\cite{Hernandez2018}, el dise\~{n}o es
\textbf{cuasiexperimental de alcance explicativo}, ya que se
manipula deliberadamente la variable independiente (arquitectura
del modelo) y se miden sus efectos sobre las variables dependientes
(mAP@0.5 y FPS) bajo condiciones controladas de oclusi\'on y
baja visibilidad.

La metodolog\'ia se basa directamente en el pipeline establecido
por Weng \& Niu~\cite{ElsYolo2025} en ELS-YOLO, adoptando su
misma configuraci\'on de entrenamiento ---optimizador, \'epocas,
batch size y dataset--- para garantizar una comparaci\'on
reproducible y justa. Sobre esa base se integran dos
contribuciones propias: un m\'odulo de atenci\'on cruzada
canal-espacial para oclusi\'on, inspirado en el mecanismo EAFF de
Shen et al.~\cite{Shen2025}, y una estrategia de augmentaci\'on
sint\'etica de oclusi\'on durante el entrenamiento.

El proceso metodol\'ogico comprende cuatro fases secuenciales:

\begin{enumerate}
  \item \textbf{Fase 1 --- Preparaci\'on del entorno y dataset:}
        descarga y verificaci\'on del dataset ExDark~\cite{Loh2019},
        conversi\'on de anotaciones a formato YOLO (.txt),
        implementaci\'on del script de augmentaci\'on sint\'etica
        de oclusi\'on, y divisi\'on del dataset en conjuntos de
        entrenamiento (70\,\%), validaci\'on (10\,\%) y prueba
        (20\,\%).

  \item \textbf{Fase 2 --- Dise\~{n}o e implementaci\'on de
        ARGOS-11:} modificaci\'on de la arquitectura YOLOv11s
        integrando el m\'odulo de atenci\'on cruzada canal-espacial
        en el cuello de la red (\textit{neck}), configuraci\'on de
        la pipeline de augmentaci\'on con Albumentations, y
        verificaci\'on de la integridad del modelo antes del
        entrenamiento.

  \item \textbf{Fase 3 --- Entrenamiento y ajuste:} entrenamiento
        del modelo ARGOS-11 siguiendo el pipeline de ELS-YOLO con
        los hiperpar\'ametros definidos en la
        Tabla~\ref{tab:hiperparametros}. Se ejecutar\'an tres
        corridas independientes para garantizar estabilidad de los
        resultados.

  \item \textbf{Fase 4 --- Evaluaci\'on y an\'alisis:} inferencia
        sobre el conjunto de prueba de ExDark, c\'alculo de m\'etricas
        (mAP@0.5, mAP@0.5:0.95, FPS), comparaci\'on directa contra
        ELS-YOLO como baseline, y an\'alisis estad\'{i}stico de los
        resultados.
\end{enumerate}

\subsection{Dise\~{n}o de la Construcci\'on de la Propuesta de
Ingenier\'ia}

La propuesta de ingenier\'ia consiste en el modelo \textbf{ARGOS-11}
(\textit{Adverse-condition Robust person detectiOn System}), una
arquitectura basada en YOLOv11s que extiende el modelo ELS-YOLO
con dos componentes adicionales para el manejo simult\'aneo de
oclusi\'on y baja visibilidad. A continuaci\'on se detallan las
etapas de construcci\'on:

\subsubsection{Arquitectura base: YOLOv11s}

Se utiliza YOLOv11s (Ultralytics, 2024) como red troncal, en su
variante peque\~{n}a (\textit{small}), por ofrecer el mejor
equilibrio entre precisi\'on y velocidad de inferencia en la
familia YOLOv11. Siguiendo a Weng \&
Niu~\cite{ElsYolo2025}, no se utilizan pesos preentrenados
(\textit{from scratch}), lo que asegura que los resultados
reflejan exclusivamente el aporte de la arquitectura y los datos,
y no de transferencia de aprendizaje.

\subsubsection{M\'odulo de atenci\'on cruzada canal-espacial
(contribuci\'on 1)}

Se integra un m\'odulo de atenci\'on cruzada canal-espacial en el
cuello (\textit{neck}) de YOLOv11s, entre el FPN y el cabezal de
detecci\'on. Este m\'odulo est\'a inspirado en el mecanismo EAFF
(\textit{Efficient Attention Feature Fusion}) de Shen et
al.~\cite{Shen2025}, adaptado para im\'agenes RGB en lugar de
pares multimodales. Su funcionamiento es el siguiente:

\begin{enumerate}
  \item \textbf{Rama de atenci\'on de canal:} aplica
        \textit{Global Average Pooling} sobre cada mapa de
        caracter\'{i}sticas para obtener un vector de pesos por
        canal, que pondera las caracter\'{i}sticas m\'as informativas
        de las regiones visibles de la persona.
  \item \textbf{Rama de atenci\'on espacial:} genera un mapa de
        activaci\'on espacial mediante convoluciones 7$\times$7
        que localiza las regiones del mapa de caracter\'{i}sticas
        con mayor probabilidad de contener siluetas parcialmente
        ocluidas.
  \item \textbf{Fusi\'on cruzada:} los mapas de las dos ramas se
        multiplican elemento a elemento (\textit{element-wise
        multiplication}) para producir un mapa de
        caracter\'{i}sticas refinado que atenua el ruido de fondo
        y enfatiza las regiones de inter\'es.
\end{enumerate}

Este mecanismo incrementa la sensibilidad del detector ante
siluetas parcialmente cubiertas sin aumentar significativamente
la complejidad computacional del modelo, lo que es
cr\'{i}tico para mantener el requisito de inferencia en tiempo
real ($>$30 FPS).

\subsubsection{Augmentaci\'on sint\'etica de oclusi\'on
(contribuci\'on 2)}

Durante el entrenamiento se aplica una estrategia de augmentaci\'on
sint\'etica de oclusi\'on sobre las im\'agenes de ExDark, inspirada
en el protocolo de ocultamiento aleatorio (\textit{random erasing})
utilizado en trabajos de oclusi\'on peatonal~\cite{Hu2025, Kim2025}.
El procedimiento es el siguiente:

\begin{enumerate}
  \item Para cada imagen de entrenamiento, con probabilidad
        $p = 0{,}5$, se selecciona aleatoriamente entre 1 y 3
        personas anotadas.
  \item Sobre el bounding box de cada persona seleccionada, se
        coloca un parche opaco rectangular de color aleatorio
        (negro, gris o con textura de fondo) que cubre entre el
        30\,\% y el 70\,\% del \'area del bounding box.
  \item La posici\'on del parche dentro del bounding box es
        aleatoria, cubriendo distintas partes del cuerpo (cabeza,
        torso, piernas) para simular diferentes tipos de oclusi\'on
        real.
  \item Las anotaciones originales se mantienen sin modificar, de
        forma que el modelo aprende a detectar la persona a partir
        de la fracci\'on visible de la silueta.
\end{enumerate}

Esta estrategia incrementa artificialmente la diversidad de los
datos de entrenamiento en t\'erminos de oclusi\'on, sin necesidad
de recopilar im\'agenes adicionales, y ha demostrado mejorar
consistentemente la robustez de los detectores ante oclusi\'on
real~\cite{Tang2026}.

\subsubsection{Configuraci\'on de entrenamiento}

La Tabla~\ref{tab:hiperparametros} detalla la configuraci\'on
completa de entrenamiento, replicando fielmente el pipeline de
ELS-YOLO~\cite{ElsYolo2025} para garantizar comparabilidad directa
de resultados.

\begin{table}[h!]
\centering
\caption{Hiperpar\'ametros de entrenamiento de ARGOS-11 (basados
en ELS-YOLO~\cite{ElsYolo2025})}
\label{tab:hiperparametros}
\renewcommand{\arraystretch}{1.3}
\begin{tabular}{|l|c|l|}
\hline
\textbf{Hiperpar\'ametro} & \textbf{Valor} & \textbf{Justificaci\'on} \\
\hline
Arquitectura base & YOLOv11s & Mismo que ELS-YOLO \\
\hline
Optimizador & SGD & Mismo que ELS-YOLO \\
\hline
Tasa de aprendizaje inicial & 0{,}01 & Mismo que ELS-YOLO \\
\hline
Momentum & 0{,}937 & Mismo que ELS-YOLO \\
\hline
Weight decay & 0{,}0005 & Mismo que ELS-YOLO \\
\hline
\'Epocas & 200 & Mismo que ELS-YOLO \\
\hline
Batch size & 16 & Mismo que ELS-YOLO \\
\hline
Pesos iniciales & Sin preentrenar & Mismo que ELS-YOLO \\
\hline
Tama\~{n}o de entrada & 640$\times$640 px & Est\'andar YOLO \\
\hline
Dataset & ExDark & Mismo que ELS-YOLO \\
\hline
Divisi\'on & 70 / 10 / 20\,\% & Train / Val / Test \\
\hline
Augmentaci\'on sint\'etica & $p = 0{,}5$, oclusi\'on 30--70\,\% & Contribuci\'on propia \\
\hline
Corridas independientes & 3 & Estabilidad estad\'{i}stica \\
\hline
\end{tabular}
\end{table}

\subsubsection{Modelo de comparaci\'on (baseline)}

El modelo de referencia es \textbf{ELS-YOLO}~\cite{ElsYolo2025},
entrenado bajo exactamente los mismos hiperpar\'ametros y dataset
que ARGOS-11. La comparaci\'on es justa porque ambos modelos
comparten la misma arquitectura base (YOLOv11s), el mismo
optimizador, las mismas \'epocas y el mismo conjunto de datos.
La \'unica diferencia es la presencia del m\'odulo de atenci\'on
cruzada y la augmentaci\'on sint\'etica de oclusi\'on en ARGOS-11.

La meta establecida es que ARGOS-11 alcance un
\textbf{mAP@0.5 $\geq$ 74{,}3\,\%} (resultado reportado por
ELS-YOLO) manteniendo una tasa de inferencia superior a 30 FPS.

\subsection{Dise\~{n}o Muestral}

\textbf{Dataset:} \textbf{ExDark}~\cite{Loh2019} (\textit{Exclusively
Dark Image Dataset}), un conjunto de datos p\'ublico y de acceso
libre compuesto por \textbf{7{,}363 im\'agenes} capturadas
exclusivamente en condiciones de baja iluminaci\'on real: interiores
oscuros, noche urbana, niebla, lluvia, luz de velas y entornos
subexpuestos. Todas las im\'agenes son \textbf{solo RGB}, sin
sensores especializados, lo que garantiza la aplicabilidad del
modelo en cualquier c\'amara convencional.

\begin{table}[h!]
\centering
\caption{Caracter\'{i}sticas del dataset ExDark}
\label{tab:exdark}
\renewcommand{\arraystretch}{1.3}
\begin{tabular}{|l|c|}
\hline
\textbf{Caracter\'{i}stica} & \textbf{Valor} \\
\hline
Total de im\'agenes & 7{,}363 \\
\hline
Clases de objetos & 12 (incluye \textit{person}) \\
\hline
Tipo de sensor & Solo RGB (c\'amaras convencionales) \\
\hline
Condiciones de iluminaci\'on & Baja (noche, interior, niebla, lluvia) \\
\hline
Formato de anotaci\'on & Bounding box (convertido a YOLO .txt) \\
\hline
Divisi\'on train / val / test & 5{,}154 / 737 / 1{,}472 im\'agenes \\
\hline
Acceso & P\'ublico y libre \\
\hline
\end{tabular}
\end{table}

\noindent\textbf{Unidad de an\'alisis:} imagen RGB de baja
iluminaci\'on que contenga al menos una persona anotada, con o sin
augmentaci\'on sint\'etica de oclusi\'on aplicada.

\bigskip
\textbf{Criterios de inclusi\'on de im\'agenes:}
\begin{itemize}
  \item Im\'agenes con al menos una persona anotada con bounding
        box.
  \item Im\'agenes capturadas en condiciones de baja iluminaci\'on
        real (todas las im\'agenes de ExDark cumplen este criterio).
  \item Im\'agenes de formato RGB est\'andar (.jpg, .png).
\end{itemize}

\textbf{Criterios de exclusi\'on de im\'agenes:}
\begin{itemize}
  \item Im\'agenes corruptas o con anotaciones inconsistentes.
  \item Im\'agenes sin ninguna persona en la escena.
\end{itemize}

\subsection{T\'ecnicas de Recolecci\'on de Datos}

Los datos experimentales provienen de dos fuentes:

\textbf{Fuente 1 --- Dataset ExDark:}
El dataset se obtendr\'a del repositorio oficial de
Loh \& Chan~\cite{Loh2019}
(\url{https://github.com/cs-chan/Exclusively-Dark-Image-Dataset}).
Las anotaciones originales se convertir\'an a formato YOLO (.txt)
mediante un script Python propio. El proceso incluye:
(1)~descarga y verificaci\'on de integridad (MD5);
(2)~conversi\'on de coordenadas de bounding box al formato
normalizado YOLO;
(3)~verificaci\'on de que todas las categor\'{i}as de persona est\'en
correctamente etiquetadas;
(4)~divisi\'on aleatoria estratificada en conjuntos de
entrenamiento (70\,\%), validaci\'on (10\,\%) y prueba (20\,\%).

\textbf{Fuente 2 --- M\'etricas de entrenamiento:}
Durante el entrenamiento se registrar\'an autom\'aticamente mediante
Ultralytics YOLO las siguientes m\'etricas por \'epoca:
p\'erdida de localizaci\'on (\textit{box loss}), p\'erdida de
clasificaci\'on (\textit{cls loss}), mAP@0.5 y mAP@0.5:0.95 en el
conjunto de validaci\'on. Estos datos se exportar\'an en formato CSV
para su posterior an\'alisis estad\'{i}stico.

\subsection{T\'ecnicas Estad\'{i}sticas para el Procesamiento de
la Informaci\'on}

Para la evaluaci\'on y comparaci\'on de ARGOS-11 contra ELS-YOLO
se emplear\'an las siguientes t\'ecnicas estad\'{i}sticas:

\begin{itemize}
  \item \textbf{Prueba \textit{t} de Student (muestras
        independientes):} para comparar la distribuci\'on de
        mAP@0.5 entre ARGOS-11 y ELS-YOLO a lo largo de las tres
        corridas independientes, con nivel de significancia
        $\alpha = 0{,}05$.

  \item \textbf{Intervalo de confianza al 95\,\%:} para reportar
        la incertidumbre en las estimaciones de mAP@0.5 y FPS,
        expresando los resultados como media $\pm$ desviaci\'on
        est\'andar sobre las tres corridas.

  \item \textbf{Curvas Precision-Recall:} para visualizar el
        comportamiento del detector en distintos umbrales de
        confianza y calcular el \'area bajo la curva (mAP).

  \item \textbf{An\'alisis de ablaci\'on:} para cuantificar el
        aporte individual de cada componente de ARGOS-11 mediante
        tres variantes de entrenamiento:
        \begin{enumerate}
          \item YOLOv11s base (= ELS-YOLO, sin modificaciones)
          \item YOLOv11s + m\'odulo de atenci\'on (sin augmentaci\'on)
          \item YOLOv11s + augmentaci\'on (sin m\'odulo)
          \item ARGOS-11 completo (ambos componentes)
        \end{enumerate}
        Esto permite aislar la contribuci\'on de cada aporte y
        determinar cu\'al tiene mayor impacto en el mAP@0.5.

  \item \textbf{Software:} Python 3.11 con Ultralytics YOLOv11,
        PyTorch, Albumentations (augmentaci\'on), pycocotools
        (m\'etricas COCO), numpy, pandas, matplotlib y
        scipy.stats. El entrenamiento se ejecutar\'a bajo tres
        alternativas de infraestructura: (A)~Google Colab Pro+,
        (B)~NVIDIA Cloud en Brev.dev con RTX Pro 6000, o
        (C)~GPU local RTX 5070 Ti de 16\,GB. La validaci\'on de
        FPS se realizar\'a sobre dispositivo embebido Jetson Nano.
\end{itemize}

\subsection{Aspectos \'Eticos}

Esta investigaci\'on es de car\'acter experimental no invasivo y
no involucra datos personales, participantes humanos ni animales.
El dataset ExDark es p\'ublico y de acceso libre bajo licencia
acad\'emica, sin restricciones de uso para investigaci\'on
cient\'{i}fica~\cite{Loh2019}. Todos los art\'{i}culos referenciados
son citados correctamente siguiendo el estilo IEEE. Los autores
declaran no tener conflicto de intereses y que el financiamiento
es exclusivamente institucional (asignatura Investigaci\'on III,
Universidad Peruana Uni\'on, Juliaca). La originalidad del trabajo
fue verificada con un \'{i}ndice de similitud del 10\,\% mediante
Turnitin, en cumplimiento de las normas de integridad acad\'emica
de la UPeU.
